\properlane{renewed-action}{The Renewal of Hands, Words, and Desires}

The new person is recognizable in the use of ordinary human powers.
Speech becomes truthful, anger is restrained, hands work to share,
and desire seeks the ways of God. Chrysostom's practical exegesis
and Bellarmine's account of grace make this more than a programme
of self-improvement. God renews a person who must really act;
the action never becomes independent of its giver. The Mass gathers
these powers in prayer and seeks their healing for the life beyond
the celebration.

\subsection{The same person, newly alive}

Paul does not ask the hearer to exchange his human nature for
another. Chrysostom insists on the identity of the person who is
renewed: the change concerns the manner of life. The command to
put on the new person would be unintelligible if the old way of
living were the person's unalterable essence. Equally, the language
of creation says that renewal is more than a new arrangement of
unchanged habits under a new name. God brings about the life that
Paul commands.\footnote{Chrysostom, \work{Ephesians} XIII, on
4:23--24.}

The combination matters for repentance. If evil conduct defined
the person beyond remedy, the command could only condemn. If a
new name were sufficient, the command would require no conversion.
Chrysostom's account allows the same person to acknowledge what
he has done and to live differently through the grace he receives.
Baptismal clothing has a lasting demand: the given garment must
remain visible in conduct.

Aquinas distinguishes several possible explanations of the
``spirit of your mind'': the Holy Spirit dwelling in the mind,
the rational spirit, or a mind made spiritual. He does not make
one grammatical phrase carry an artificially simple psychology.
His exposition nevertheless converges on the concrete renewal
of holiness in the heart, truth in the mouth, and justice in
action. Christ is the beginning of newness where Adam names
oldness.\footnote{Thomas Aquinas, \work{Super Ephesios} 4,
lectio 7; Dessain 1857, vol.~II, p.~329. Paraphrase of the
inspected Latin text.}

The Collect's body and mind belong within this renewal.
Schuster reads their ordering toward God as the freedom of the
whole person. Embodiment is not the obstacle to holiness. The
problem is the disorder that keeps human powers from serving
their true end. Hands, lips, attention, and desires can all be
freed for service. The Introit names the first such action:
incline the ear. Renewal begins with receiving a word that
comes from God, not with the announcement of an independent
project of moral mastery.

\subsection{Truth protects the body}

Chrysostom's image of the eye and foot uncovers the social
structure of lying. One member supplies another with knowledge
it needs in order to move. If the eye conceals danger, the foot
falls; if the foot misreports the ground, the eye's body suffers
too. A lie corrupts shared life at the point of dependence.
The fact that the liar has gained an immediate advantage cannot
make the act harmless to the body he belongs to.

Paul's reason is stronger than a rule against impolite speech.
``Members one of another'' means that another's safety and
good belong within one's own concern. Truth enables the
neighbour's action; falsehood manipulates it. A renewed person
therefore asks not merely whether a sentence can survive a
technical defence, but whether the neighbour receives the
truth on which he must rely. This is an application of the
body's mutual service, not a licence for indiscriminate
disclosure of everything one happens to know.

The psalm surrounding the Gradual makes the mouth itself an
object of prayer. Bellarmine explains its watch and door as
the help needed to know when to speak and to have the strength
to speak or remain silent accordingly. The understanding and
will both need assistance. He then adds that cooperation with
grace is itself a gift; it supplies no reason for boasting.
\footnote{Chrysostom, \work{Ephesians} XIV, on 4:25;
Bellarmine, \work{Psalms} 140:3. Psalm 140:3 supplies context
for the appointed verse 2.}

The Alleluia gives this guarded speech an affirmative task.
It proclaims God's deeds among the nations. Truthfulness
within the community and witness beyond it are different
acts of the same renewed life. A tongue carefully restrained
from falsehood has not thereby reached its only purpose.
It can praise and invite. Bellarmine's account of love wishing
the beloved to be known makes proclamation a form of
generosity, while his insistence on help even for praise
keeps the speaker dependent on God.

\subsection{Anger must not become a dwelling place}

The Epistle does not leave anger as a merely interior matter.
It can become the breach through which a common life is
divided. Chrysostom follows its development from an injury
to repeated inward rehearsal, and from rehearsal to distorted
judgment. When enmity takes the judge's seat, a doubtful
word sounds hostile, an innocent gesture becomes an insult,
and a slander finds a ready advocate. The danger is not only
that a person feels too strongly. He begins to perceive
his neighbour through the settled aim of opposing him.
\footnote{Chrysostom, \work{Ephesians} XIV, on 4:26--27.}

Sunset gives that process a boundary. Chrysostom's account
of night shows why promptness matters: fewer outward
distractions leave more room for feeding resentment.
Reconciliation is not something to consider only after
anger has matured into a complete case against the other
person. Its urgency belongs to the struggle against that
maturation. One cannot assume that waiting leaves the
heart neutral.

The command does not require a fiction that every wrong
has been repaired by evening, nor that emotional agitation
can be ended by consulting a clock. A wrong may still
require truthful correction. The prohibition concerns
giving anger room to become a continuing occupation of
the heart. A person can refuse revenge, stop nursing an
accusation, and seek peace before all the circumstances
are resolved. Chrysostom directs hostility against the
real enemy of the body's unity rather than against a
fellow member.

The Offertory helps sustain this refusal. God's saving
hand, not the worshipper's resentment, is the final source
of life. The prayer acknowledges enemies without making
hatred the means of salvation. Augustine reads secure
life beyond their reach and rejects proud reliance on
one's own works. The person seeking reconciliation can
therefore act without pretending to control the outcome.
\footnote{Augustine, \work{Enarrationes} 137.10--13;
Bellarmine, \work{Psalms} 137:7--8.}

\subsection{Work finds its purpose in giving}

Paul does not say merely that the former thief must become
economically successful. Honest work is ordered toward
having something to give to the person in need. The
change is relational: hands once used to take are now
used to supply another. Chrysostom insists on this
positive movement. Stopping the theft matters, but the
new life is not exhausted by an absence of theft.
\footnote{Chrysostom, \work{Ephesians} XIV, on 4:28.}

The command also distinguishes good work from effort
considered in itself. Diligence can be directed toward
harm. The thief's labour is real labour, but its object
and purpose are wrong. Nor is work purified simply by
its power to enrich the worker. Paul's final clause
makes the neighbour's need part of its purpose. A
renewed economy of action asks both whether the work
is good and whether its fruits can become a benefit
outside the self.

Gregory's explanation of the bound hands makes a
searching comparison with this Epistle. The guest's
hands and feet are bound in judgment; Gregory says
that culpable refusal has already bound them to
inaction in life. Hands that give nothing to those
in need and feet that neglect the sick are withheld
from the good they could do. The judicial binding
reveals the seriousness of that voluntary withholding.
\footnote{Gregory, \work{Gospel Homilies} 38.13.}

His charge concerns refusal, not physical inability.
Dependence, disability, poverty, or lack of paid work
is not the absent garment. Paul names the person
in need as the one to whom help is due, not as an
unsuccessful worker who has failed to justify his
place at the table. Gregory's moral reading examines
what a person withholds when he can help. It does
not turn human strength into the measure of worth.

The Gospel also keeps labour from becoming a
self-contained source of meaning. The invitation
is to the Son's feast; those who prefer farm or
merchandise can refuse it while remaining busy.
Work becomes a good within a relationship with
God, not an adequate replacement for that
relationship. The hands that share belong to
a guest, someone first called to receive.

\subsection{Raised hands belong to Christ's offering}

The Gradual brings action into prayer. Its lifted
hands are the bodily expression of a petition for
acceptance. Augustine's evening sacrifice is the
voluntary Passion of Christ; the Church's prayer
belongs to the body whose Head has offered himself.
The worker's generosity is therefore not the
independent ground on which he forces entrance
before God. His good action and his prayer
participate in a life given by Christ.

Bellarmine develops the same dependence through
prayer's intention. Incense must rise toward God;
the wish to attract spectators draws it aside.
A visible gesture can conceal an inward aim
at admiration. Prayer needs faith, hope, love,
humility, and Christ's priestly mediation.
The request for prayer to be directed is itself
an admission that the worshipper cannot assume
these conditions are already perfectly present.
\footnote{Augustine, \work{Enarrationes} 140.2--4;
Bellarmine, \work{Psalms} 140:1--2.}

This examination of intention returns to honest
work. The same person can seek admiration while
giving, or praise while praying. External action
matters, but its purpose matters too. The secret
of generosity cannot be resolved by how publicly
the hands are displayed. Charity directs the act
toward God and the neighbour; vanity directs it
back toward the actor's reputation.

The Secret asks that the offering be saving for
the participants. Schuster places this request
at the point where the gift's power and the
recipient's disposition meet. Human labour's
gift is not simply identical with the Eucharistic
offering, and neither is another independent
Passion. The Eucharist is God's saving gift
through which the participant's life is
transformed. The worker comes to receive
the grace by which his giving can become
charity rather than an achievement possessed
against God or neighbour.

\subsection{Desire itself needs to be healed}

Bellarmine's treatment of Psalm 118 prevents
the language of action from becoming
self-justification. God's command is real,
and initial justification is not the reward
for unaided observance. Grace precedes the
obedience through which the justified person
grows. At verse 5 the wish for directed
ways becomes a request for the assistance
without which the command will not be
fulfilled. At verse 8 even a firm resolution
is joined to a plea not to be forsaken.
\footnote{Bellarmine, \work{Psalms} 118:4--8,
especially 5 and 8. Verse 8 is contextual
reception; the Communion appoints 4--5.}

Augustine likewise hears a prayer in the
wish. The Communion does not oppose command
to grace, as though one must choose either
serious obligation or real dependence. It
utters them in sequence. God commands;
the worshipper asks that his life be
directed to obey. Desire can acknowledge
both the goodness of the command and the
weakness of the one who loves it.

The Postcommunion carries the argument
beneath individual acts to the inclination
from which they arise. Its medicine frees
from perversity and produces adherence.
Schuster interprets the Eucharistic grace
as healing the desires of a person whose
spiritual life needs transformation.
Moral conversion is therefore more than
performing a correct action while wishing
with one's whole heart to do the opposite.
The prayer asks for the freedom to love
and continue in the good.

This healing neither renders action
unnecessary nor turns repeated struggle
into proof that grace is absent. The
petition's continuing force is shown by
its request for \textit{always}: a life
must be sustained, not merely begun.
The next truthful word, the interruption
of resentment, and the gift to need are
particular places where that adherence
can become real. Each is small enough
to do; none is separate from the life
God is renewing.

The final feast is the horizon of this
renewal. Human powers are made ready for
communion with the Son, not trained
merely for more efficient activity.
Augustine's Passion reading of the
evening sacrifice and Schuster's hope
of resurrection morning place present
labour within an end greater than its
visible success. Judgment gives the
present act seriousness; resurrection
gives it hope.\footnote{Schuster,
\work{Sacramentary} III, pp.~172--174;
Augustine, \work{Enarrationes} 140.3.}

\clearpage
\begin{fourSenses}
\item[Literal.] Paul addresses the renewed community's speech, anger,
and labour. The chants ask for hearing, accepted prayer, proclamation,
preservation, and directed ways; the orations seek unhindered service
and medicinal change. The royal invitation makes the response consequential.
\item[Allegorical.] Christ is the source of the new person's life. His
Passion is the evening sacrifice in which his body prays. The Bridegroom's
charity gives renewed action its form and the feast its purpose.
\item[Moral.] Tell the truth on which another depends, interrupt the
cultivation of resentment, and work to share with need. Ask for God's
help in intention as well as action. Culpable refusal differs from
inability; neither poverty nor weakness cancels the invitation.
\item[Anagogical.] Present renewal tends toward perfected human life
in the Son's final feast. Trial and labour are not the last word:
Christ's evening sacrifice opens upon resurrection. Final judgment
refuses endless postponement, while healing grace sustains perseverance.
\end{fourSenses}
